%=====================================================================
% Front matter: title page, certificate, declaration, acknowledgement,
% abstract, contents, list of figures, list of tables.
%=====================================================================
\begin{titlepage}
\begin{tikzpicture}[remember picture,overlay]
  \draw[line width=2.5pt] ([xshift=0.6in,yshift=-0.6in]current page.north west)
    rectangle ([xshift=-0.6in,yshift=0.6in]current page.south east);
\end{tikzpicture}
\centering
\thispagestyle{empty}
\begin{spacing}{1.0}
\vspace*{0.02in}

\includegraphics[width=1.6in]{figures/pes_logo.png}\par
\vspace{0.2in}
{\large\itshape Dissertation on\par}
\vspace{0.5em}
{\Large\bfseries ``Tinker RL Lab:\\[0.35em]
Post-Training of Large Language Models''\par}
\vspace{0.4em}
{\normalsize A multi-framework training harness and studies of GRPO signal starvation\par}

\vfill

{\normalsize\itshape Submitted in partial fulfilment of the requirements for the award of the degree of\par}
\vspace{0.6em}
{\large\bfseries Master of Technology\par}
\vspace{0.2em}
{\normalsize\bfseries in\par}
\vspace{0.2em}
{\large\bfseries Data Science and Machine Learning\par}
\vspace{0.2em}
{\large\bfseries UE20CS972 -- Project Phase -- 2\par}

\vfill

{\normalsize\itshape Submitted by:\par}
\vspace{0.3em}
{\large\bfseries Arvind C R \qquad PES2PGE24DS140\par}

\vspace{0.5em}
{\normalsize\itshape Under the guidance of\par}
\vspace{0.2em}
{\large\bfseries Ramesh Prakash Guledgudd\par}

\vspace{0.6em}
{\normalsize\bfseries\color{pesmaroon} May 2026 -- October 2026\par}

\vfill

{\normalsize\bfseries DEPARTMENT OF COMPUTER SCIENCE AND ENGINEERING\par}
{\normalsize FACULTY OF ENGINEERING\par}
{\large\bfseries PES UNIVERSITY\par}
{\footnotesize (Established under Karnataka Act No.\ 16 of 2013)\\
Electronic City, Hosur Road, Bengaluru -- 560 100, Karnataka, India\par}
\end{spacing}
\end{titlepage}

%---------------------------------------------------------------------
\begin{titlepage}
\centering
\thispagestyle{empty}
\vspace*{0.1in}
{\Large\bfseries PES UNIVERSITY\par}
\vspace{0.2in}
{\large\bfseries FACULTY OF ENGINEERING\par}
{\normalsize DEPARTMENT OF COMPUTER SCIENCE AND ENGINEERING\par}
\vspace{0.35in}
{\Large\bfseries CERTIFICATE\par}
\vspace{0.5in}
\begin{spacing}{1.4}
\noindent This is to certify that the project report entitled
\textbf{``Tinker RL Lab: Post-Training of Large Language Models --- A
multi-framework training harness and studies of GRPO signal starvation''} is a
bonafide record of the project work carried out by
\textbf{Arvind C R} (SRN: \textbf{PES2PGE24DS140}), in partial fulfilment of
the requirements for the award of the degree of Master of Technology in
Data Science and Machine Learning of PES University, during the period
May 2026 -- October 2026. It is certified that all corrections/suggestions
indicated for internal assessment have been incorporated in the report.

\vspace{0.3in}
The candidate's individual contribution, the inherited Semester-3 group
infrastructure, and the boundary between them are stated in the Declaration
and in Chapter 1 \S1.6.
\end{spacing}
\vspace{0.4in}
\noindent\begin{tabular}{@{}p{2.4in}p{2.4in}@{}}
\dotfill & \dotfill \\
\textbf{Ramesh Prakash Guledgudd} & \textbf{Head of the Department} \\
Project Guide & Department of CSE \\
\end{tabular}
\vspace{0.3in}
\noindent\begin{tabular}{@{}p{2.4in}@{}}
\dotfill \\
\textbf{External Examiner} \\
\end{tabular}
\vfill
{\normalsize October 2026\par}
\end{titlepage}


%---------------------------------------------------------------------
\chapter*{Declaration}
\addcontentsline{toc}{chapter}{Declaration}
\thispagestyle{plain}

I declare that this project report is my own work. Where the work of others has
been used, it is cited and acknowledged. The matter embodied in this report has
not been submitted, in part or in full, to this or any other university or
institution for the award of any other degree or diploma.

The multi-framework benchmark infrastructure, the baseline framework
integrations, and the initial GRPO/PPO/DPO experimentation were produced by the
Semester-3 project group, of which I was a member, and are inherited here as
the shared foundation. The research direction, the experimental expansion, the
statistical analyses, the P1--P8 internal reports, the reporting and registry
artefacts, the adaptive group-size controller, the applied fraud study, and the
E1--E14 evaluation campaign reported in this document are my
individual contribution, carried out under the guidance of my project guide.

Quantitative results are tied to named evidence. Appendix B distinguishes
recovered primary campaign files, inspectable artefacts and the remaining
limits on underlying payload verification. Incomplete, blocked and negative
results are reported rather than omitted. The boundary between inherited and
individual work is stated precisely in \texttt{platform\_hybrid/sem 4
work/PROVENANCE.md}.

\vspace{0.2in}
\noindent\begin{tabular}{@{}p{2.6in}@{}}
\includegraphics[width=2.2in]{figures/signature_arvind.png} \\
\dotfill \\
\textbf{Arvind C R} \\
SRN: PES2PGE24DS140 \\
\end{tabular}

%---------------------------------------------------------------------
\chapter*{Acknowledgement}
\addcontentsline{toc}{chapter}{Acknowledgement}
\thispagestyle{plain}

I thank my project guide, \textbf{Ramesh Prakash Guledgudd}, for direction and
sustained feedback throughout this work, and for the freedom to pursue the
conservative and negative results that form a significant part of it.

I thank the Semester-3 project group for the shared multi-framework benchmark
and the baseline integrations on which the Semester-4 studies build. I thank
the maintainers of the benchmark suites used in the E1--E14 campaign, whose
native evaluators made an independently checkable evaluation possible, and the
platforms --- Tinker, Modal, and Google Cloud --- on which the experiments ran.

Finally, I thank the reviewers whose questions at earlier reviews shaped the
reporting discipline this report depends on: the insistence on stating what a
number does \emph{not} establish has improved the work materially.

\vspace{0.3in}
\noindent\textbf{Arvind C R}

%---------------------------------------------------------------------
\chapter*{Abstract}
\addcontentsline{toc}{chapter}{Abstract}
\thispagestyle{plain}

This report studies group-relative reinforcement-learning post-training of
large language models through an explicit record of the training and
evaluation stack, at short horizons (30--50 steps), with verifiable rewards,
on small and mid-scale models plus managed-API anchors. Controlled comparisons seek to isolate the factor under
test; remaining differences and missing evidence limit attribution.
It makes four contributions. Terminology used in every headline: a study is
labelled \emph{clipped-GRPO} only where a clipped ratio surrogate is verified in
the runner, \emph{REINFORCE group-mean} where the inspected loop centres
rewards within the group with no clipping (with or without an unclipped
importance ratio), and \emph{unknown (managed or unrecovered)} where the
executed loss cannot be inspected (Chapter~1 terminology box; Table~4.2).
Several best-audited stacks below (the group-size sweep, the four-method
controller trace and the custom Tinker loops) are therefore REINFORCE
group-mean, not GRPO as published.

\textbf{(1) ZVF instrumentation and a loss-form audit (P2 demoted to
instrumentation).} The Zero-Variance Fraction (ZVF) is logged per step and
recomputed from stored reward tensors, with a closed-form Bernoulli null.
That recomputation is an independent check on stored tensors; it is not
prognostic validation. ZVF's predictive validity is unestablished: the
historical pooled correlations are quarantined as withdrawn mixed-pool
analyses (17 finite cells, 9 synthetic, with mixed outcomes including training
reward; the 52-row AUROC panel's 41 safe rows are synthetic, with 7 imputed
collapse anchors). Reading the runners' code distinguishes custom REINFORCE
group-mean loops from clipped surrogates, nonzero-KL configuration records and
unknown managed settings.

\textbf{(2) Stack-conditioned reporting (schema proposal).} A minimum
reporting standard and a machine-readable stack registry with a
\texttt{stackdiff} tool that compares declared fields and reports
\emph{declared-difference codes} (R0--R5), not a calibrated probability of a
label flip and not an externally validated badge. Consistency of the registry
on its own seed entries is a closed-system check, not external validation.

\textbf{(3) Calibrated negative results.} No reliable scaling law is recovered
from reward curves across twelve fitted anchors (loss form unknown). The
saturated same-stack clipped-GRPO versus PPO contrast (0.992 against 0.990; exact sign-flip $p=0.625$, paired-$t$
$p=0.374$, five seeds, paired by seed not by item, with asymmetric
train--evaluation overlap) and the saturated REINFORCE group-mean group-size
sweep are
\emph{pilot/feasibility evidence only}: at ceiling they cannot separate no
effect from no headroom, and neither establishes equivalence or an optimum.
The unsaturated PPO contrast (clipped-GRPO) is demoted accordingly: the minimal
untuned-critic PPO arm (linear value head from zero, no warm-up, no GAE, no
tuning, seeing $8\times$ as many distinct prompts per step as GRPO $G=8$,
completions 190--196 tokens under a 200 cap) shows no detectable held-out
gain, while the GRPO arms gained; the post-training gap is $+0.050$
(paired-$t$ 95\% CI $[0.015, 0.085]$) with exact sign-flip $p=0.0625$ first
(the $n=5$ floor, not below 0.05) and paired-$t$ $p=0.016$ second. This is
not evidence about PPO versus GRPO as algorithms. $G = 2$ and $G = 8$ show no
detectable difference, without establishing equivalence. An adaptive
group-size controller derived from ZVF is \emph{unevaluated as a training
policy}: its plug-in ``zero rescue'' is imposed by the estimator
($\hat p\in\{0,1\}$ forces $h_G=1$), and 92\% of its escalations target
observed all-correct groups, without proving latent certainty or impossibility
of recovery with fresh rollouts. No matched-budget prospective bakeoff was run.

\textbf{(4) An evaluation-governance protocol (stored-evidence, not
replayed; actor loss unknown).} One lightly trained actor
(Qwen3.6-35B-A3B with a LoRA adapter, 40 training steps; executed loss
unknown (unrecovered dirty source), adapter revision withheld, exact consumed
training rows unavailable) is evaluated on
the E1--E14 campaign, with completed, partial and blocked scopes kept separate.
All actor-dependent campaign figures are \emph{stored-evidence results, not
replayed}, and are end-to-end \emph{harness scores}, not capability measures:
for example 593/713 E1 patches were corrupt for \texttt{git apply}, E11 has
150 extraction failures under a 1,024-token thinking-on cap, and E8 has
1,259/1,967 unparsed responses. Evaluated scopes use named native graders with
the available version or revision records. A bounded 8-gram source-envelope
check over 3 of 14 suites does not establish global or semantic
decontamination. Receipts establish file consistency, not that the described
training produced the evaluated weights. The public-checkout number audit is
INCOMPLETE: of 970 registered occurrences, 828 PASS, 22 are CONTEXT-only, and
120 are ARTIFACT\_MISSING (33 of the 828 are arithmetic from printed counts and 501 are token-presence
checks against the project's own source files, i.e.\ closed-system consistency
rather than independent recomputation);
affected figures are therefore reported, not publicly verifiable from this
package alone. Receipts,
a no-pooling rule and a per-lane serving audit expose a serving defect in a
zero-scoring run; the defect's contribution to that score is unquantified.
Complete scopes include VerilogEval harness score
($129/312 = 41.35\%$), LAB-Bench harness score ($450/1967 = 0.2288$) and
AgentDojo benign-utility harness score
($88/97$); six replacement scopes completed on 27 September 2026
(Chapter~8, Table~8.C); seven original contracts remain externally blocked. A
paired trained-versus-base rerun on 5--151 identical items per lane finds no
significant difference in any lane (inconclusive at those sizes). No
claim of campaign-level improvement over the base model is made.


\textbf{C1 diagnostic extension (parser/reference audit case study, not a
recovery-rate estimate; assistant-assisted, without independent human
validation).} Benchmark reference errors, ambiguous wording and equivalent units limit the interpretation of a new frozen-policy study on public GSM8K train. The first 64 of 87 eligible questions from 4,096 screened questions received 24 fresh responses after eight valid, uncapped parser-zero responses. The frozen endpoint is 11/64 (17.19\%; Wilson 95\% interval 9.88--28.21\%). Assistant-assisted reconciliation corroborates three numerical witnesses on the same denominator, while preserving convention-dependent and ambiguous cases and one wrong-reference match. This is conditional diagnostic evidence, not a population estimate, reasoning validation or training gain. A separate, undeployed parser candidate preserves 2,025 previously valid parse objects and admits four unambiguous extractions in a predefined future subset of 2,048 outputs. It does not alter C1's frozen parser-v2 score.

\noindent\textbf{Claim map (19 October 2026).} Examinable contribution:
ZVF as recomputable \emph{instrumentation} (Ch.~6 \S6.3.1--6.3.3); the
loss-form audit and terminology box (Ch.~1, Table~4.2); stack-conditioned
reporting as a \emph{schema proposal} with declared-difference codes
(Ch.~7 \S\S7.1--7.2); calibrated pilot/null contrasts (Ch.~6); and the
evaluation-governance protocol over stored-evidence harness scores (Ch.~8).
Archival/withdrawn, not examinable as positive findings: mixed-pool ZVF
prognostic correlations and AUROCs (\S6.3.4, quarantined), the adaptive-G
controller as a training policy (\S7.3, unevaluated), saturated contrasts
beyond pilot status, and Appendix~F's off-thesis fraud side-study.

\vspace{0.25in}
\noindent\textbf{Keywords:} GRPO, RLVR, zero-variance fraction, signal
starvation, stack-conditioned results, benchmark evaluation, reproducibility,
evaluation governance, attribution.

%---------------------------------------------------------------------
\tableofcontents
\listoffigures
\listoftables
